%% Anexo B — Acta de constitución de la consultora Legasoft
%% Fragmento: no se compila solo. Lo incluye docs/entrega/entrega-fase1.tex.

\section{Anexo B. Acta de constitución de la consultora}

\begin{nota}[Objeto de este documento]
La presente acta autoriza formalmente el inicio de actividades de la consultora
Legasoft durante el \lgSemestre. No corresponde al acta de constitución de un
proyecto con cliente, que se emitirá por separado con el código
\texttt{LGS-PRY-001} cuando se seleccione la organización cliente.
\end{nota}

En la ciudad de La Paz, Estado Plurinacional de Bolivia, a los doce días del mes
de agosto de 2026, reunidos los integrantes que suscriben al final del presente
documento, en el marco de la asignatura \lgMateria\ de la carrera de
\lgCarrera\ de la Universidad Católica Boliviana «San Pablo», acuerdan
constituir la consultora \textbf{Legasoft} y dejar constancia de lo siguiente.

% ---------------------------------------------------------------------------
\subsection{Primera: constitución y denominación}

Se constituye una consultora simulada de transformación digital e ingeniería de
software que girará bajo la denominación de \textbf{Legasoft}. La consultora
tiene carácter académico y se crea con fines formativos; no adquiere personería
jurídica ni realiza actividad comercial. Sus miembros asumen, sin embargo, los
compromisos aquí descritos con el mismo rigor que exigiría una empresa real.

% ---------------------------------------------------------------------------
\subsection{Segunda: objeto}

Legasoft tiene por objeto prestar servicios de consultoría y desarrollo de
software conforme a las líneas de servicio establecidas en el documento
\texttt{LGS-INST-001}: consultoría y transformación digital; ingeniería y
desarrollo de software; backend, APIs y datos; experiencia de usuario y
frontend; Cloud y DevOps; e inteligencia artificial y analítica.

Durante el periodo de vigencia de esta acta, el objeto se concreta en atender a
una organización cliente, levantar y documentar sus requisitos, diseñar y
construir una solución de software y verificar su calidad, produciendo la
documentación técnica correspondiente.

% ---------------------------------------------------------------------------
\subsection{Tercera: domicilio y ámbito de operación}

La consultora fija su domicilio académico en la Universidad Católica Boliviana
«San Pablo», sede La Paz. Su ámbito de operación es principalmente remoto y
asincrónico, y se articula a través del repositorio institucional
\href{\lgRepo}{\lgRepoCorto}, del tablero de seguimiento en Jira y de reuniones
presenciales o virtuales según se convenga con el cliente.

% ---------------------------------------------------------------------------
\subsection{Cuarta: socios fundadores y cargos}

Son socios fundadores de Legasoft, con los cargos que se detallan, las siguientes
personas:

\begin{center}
\begin{tabularx}{\linewidth}{@{}L{4.2cm} L{4.4cm} Y@{}}
\toprule
\sffamily\bfseries\color{lgPrimario}Socio fundador &
\sffamily\bfseries\color{lgPrimario}Cargo &
\sffamily\bfseries\color{lgPrimario}Ámbito de responsabilidad \\
\midrule
\lgDirector   & \lgDirectorRol   & Relación con el cliente, requisitos, cronograma y seguimiento. \\
\lgArquitecto & \lgArquitectoRol & Arquitectura, modelo de datos, APIs y lógica de negocio. \\
\lgQA         & \lgQARol         & Interfaz de usuario, aseguramiento de calidad y pruebas. \\
\bottomrule
\end{tabularx}
\captionof{table}{Socios fundadores y cargos asignados.}
\end{center}

Los cargos son de asignación principal y no excluyente: los tres socios
participan de la planificación, las revisiones y las decisiones de alcance,
conforme a la estructura descrita en \texttt{LGS-INST-001}.

% ---------------------------------------------------------------------------
\subsection{Quinta: aportes de los socios}

El aporte de los socios a la consultora es de trabajo profesional. Cada socio se
compromete con la dedicación horaria que se indica, valorizada a la tarifa
interna definida por la consultora para el periodo:

\begin{center}
\begin{tabularx}{\linewidth}{@{}Y C{2.3cm} C{2.3cm} C{2.1cm} C{2.4cm}@{}}
\toprule
\sffamily\bfseries\color{lgPrimario}Socio &
\sffamily\bfseries\color{lgPrimario}Horas por semana &
\sffamily\bfseries\color{lgPrimario}Horas del semestre &
\sffamily\bfseries\color{lgPrimario}Tarifa (Bs/h) &
\sffamily\bfseries\color{lgPrimario}Aporte valorizado \\
\midrule
\lgDirector   & 12 & 192 & 50 & Bs 9.600 \\
\lgArquitecto & 12 & 192 & 50 & Bs 9.600 \\
\lgQA         & 12 & 192 & 50 & Bs 9.600 \\
\midrule
\textbf{Total} & \textbf{36} & \textbf{576} & — & \textbf{Bs 28.800} \\
\bottomrule
\end{tabularx}
\captionof{table}{Aportes de trabajo comprometidos para el \lgSemestre.}
\end{center}

El cálculo considera \lgSemanas\ semanas efectivas de trabajo dentro del periodo
de \lgPeriodo. La tarifa es una valorización interna de referencia, destinada a
estimar el coste del esfuerzo y a dimensionar propuestas; no implica
remuneración entre los socios ni obligación de pago alguna. El detalle de la
disponibilidad de cada socio consta en el Anexo C.

% ---------------------------------------------------------------------------
\subsection{Sexta: autorización de inicio de actividades}

\begin{nota}[Cláusula de autorización]
Los socios fundadores \textbf{autorizan formalmente el inicio de actividades de
la consultora Legasoft} a partir del 12 de agosto de 2026 y por todo el
\lgSemestre, comprendido entre \lgPeriodo. Con esta autorización quedan
habilitados la asignación de los recursos comprometidos en la cláusula quinta, el
uso del repositorio institucional, la firma de compromisos con la organización
cliente que se seleccione y la emisión de los entregables previstos.
\end{nota}

% ---------------------------------------------------------------------------
\subsection{Séptima: órgano de decisión y quórum}

Las decisiones de la consultora se adoptan en reunión de socios. Se establece:

\begin{itemize}
  \item Las decisiones ordinarias —planificación semanal, asignación de tareas,
        aceptación de entregables internos— se toman por acuerdo de los presentes,
        con un mínimo de dos socios.
  \item Las decisiones que afecten el alcance comprometido con el cliente, la
        arquitectura de la solución o las fechas de entrega requieren la
        participación de los tres socios.
  \item En caso de empate persistente, decide el \lgDirectorRol, quien deberá
        registrar el fundamento de su decisión en la minuta correspondiente.
  \item Toda decisión relevante se documenta: las funcionales en minuta, las
        arquitectónicas como decisión registrada (ADR) en
        \ruta{docs/arquitectura/}.
\end{itemize}

% ---------------------------------------------------------------------------
\subsection{Octava: compromisos de los socios}

Cada socio se obliga a:

\begin{enumerate}
  \item Cumplir la dedicación horaria comprometida y comunicar con anticipación
        cualquier indisponibilidad que afecte al cronograma.
  \item Entregar los productos de trabajo bajo su responsabilidad en las fechas
        acordadas, con la calidad y la documentación exigidas.
  \item Registrar su trabajo en el repositorio institucional y en el tablero de
        seguimiento, de modo que el avance sea verificable en cualquier momento.
  \item Someter sus entregables a revisión de al menos un par antes de
        integrarlos a la rama principal.
  \item Comunicar oportunamente riesgos, bloqueos y errores propios, en
        aplicación del valor de transparencia.
  \item Guardar reserva sobre la información que la organización cliente
        entregue, y no utilizarla para fines ajenos al proyecto.
\end{enumerate}

% ---------------------------------------------------------------------------
\subsection{Novena: régimen de reuniones}

La consultora sostendrá una reunión interna de coordinación por semana y las
reuniones con el cliente que el proyecto requiera. Toda reunión con la
organización cliente y toda reunión interna en la que se adopte una decisión
relevante deja minuta según la plantilla \texttt{LGS-MIN-NNN}, bajo
responsabilidad del \lgDirectorRol.

% ---------------------------------------------------------------------------
\subsection{Décima: normas de referencia}

La consultora adopta como marco de referencia las normas declaradas en
\texttt{LGS-INST-001}: ISO 9001, ISO/IEC 27001, ISO/IEC 12207, ISO/IEC 25010,
ISO/IEC/IEEE 29148 e ISO/IEC 20000-1. Su aplicación es progresiva y no constituye
declaración de conformidad ni certificación obtenida.

% ---------------------------------------------------------------------------
\subsection{Undécima: vigencia, modificación y cierre}

La presente acta entra en vigor en la fecha de su suscripción y permanece vigente
durante todo el \lgSemestre. Cualquier modificación de sus cláusulas —en
particular las relativas a cargos, aportes o alcance— requiere el acuerdo de los
tres socios y se formaliza mediante una nueva versión del documento, registrada
en su tabla de control de versiones y suscrita nuevamente.

Al término del periodo, la consultora cerrará sus actividades con la entrega
final al cliente, el archivo de la documentación en el repositorio institucional
y un informe de lecciones aprendidas que quedará como insumo para futuros
proyectos.

% ---------------------------------------------------------------------------
\subsection{Duodécima: aceptación}

Los socios declaran conocer y aceptar el contenido íntegro de esta acta, así como
la constitución y estructura organizacional establecidas en \texttt{LGS-INST-001},
y en señal de conformidad suscriben el presente documento.

\vspace{1.6cm}

\begin{center}
\begin{tabularx}{\linewidth}{@{}Y Y Y@{}}
\centering\rule{4.4cm}{0.5pt} &
\centering\rule{4.4cm}{0.5pt} &
\centering\arraybackslash\rule{4.4cm}{0.5pt} \\[-2pt]
\centering\sffamily\small\lgDirector &
\centering\sffamily\small\lgArquitecto &
\centering\arraybackslash\sffamily\small\lgQA \\[-4pt]
\centering\scriptsize\color{lgGris}\lgDirectorRol &
\centering\scriptsize\color{lgGris}\lgArquitectoRol &
\centering\arraybackslash\scriptsize\color{lgGris}\lgQARol \\
\end{tabularx}
\end{center}

\vspace{0.8cm}

\begin{center}
{\sffamily\small\color{lgGris}La Paz, 12 de agosto de 2026}
\end{center}
